\section{One global frame for all the local comparisons}
\label{sec:standard-global-frame}

\status{SF1--SF7 prove the common-frame and companion-lattice arguments,
including their specialization to the actual positive spaces and carried
base lattices. They replace independent local basis choices by base changes
of a single global matrix. The complete standard-column objects are still
supplied with their stated carrier properties; their arithmetic existence
is not inferred from these linear-algebra calculations.}

\paragraph{The mathematical issue.}
LI5 constructed a frame separately at every local place. The comparison
theorem needs more: one matrix over the original CM field must generate
every such frame. The source achieves this using a common self-dual lattice
and an explicit companion identity. The argument is elementary once the
actual base lattices are identified.

Write $E$ for the CM field, $\mathcal O=\mathcal O_E$, and
$\bar x$ for its CM involution. For a matrix $M$, put
$M^*=\bar M^{\mathsf T}$. In coordinates let
\[
 h_J(x,y)=\bar x^{\mathsf T}Jy,\qquad
 \Lambda^\#=\{y:h_J(x,y)\in\mathcal O\text{ for every }x\in\Lambda\}.
                                                               \tag{SF.1}
\]
The dual in (SF.1) is an $\mathcal O$-module: addition and multiplication
in the second argument preserve the integrality condition. Thus the
source's $\mathcal O$-span of \texttt{dualSet} is exactly this module.
It is a Hermitian dual, with no trace or different ideal.

\paragraph{SF1: the global matrix and the exact integer coordinates.}
Let $V$ be an $n$-dimensional $E$-space, $c:V\simeq E^n$ a coordinate
isomorphism, and $b_1,\ldots,b_n$ an orthonormal $E$-basis for $h$.
Assume $h(v,w)=h_J(c(v),c(w))$ and
$L_0=\bigoplus_j\mathcal O b_j$. Define
\[
 B=[\,c(b_1)\ \cdots\ c(b_n)\,]\in\mathrm{GL}_n(E).
\]
Then
\[
 \boxed{\ B^*JB=I,\qquad c(L_0)=B\mathcal O^n,\qquad
 x\in c(L_0)\Longleftrightarrow B^{-1}x\in\mathcal O^n.\ }       \tag{SF.2}
\]
Indeed the $(i,j)$ entry of $B^*JB$ is $h(b_i,b_j)$, while
$B^{-1}c(v)$ is the unique coefficient vector of $v$ in the basis $b$.
This proves both assertions. In the source, the first-level equality
$\operatorname{rank}(X_0)=R.k$ only reindexes this basis: apply the same
permutation to rows, columns and coefficient vectors. It neither changes
the Gram matrix nor introduces a denominator. This is the full argument
for \texttt{stdB\_core} and \texttt{exists\_stdB\_of\_tied}.

\paragraph{SF2: the actual spaces already have standard coordinates.}
For the supplied true-column record $I$ with \texttt{I.Ok}, its
\texttt{spaces} field identifies every $X_i$ with the positive space
\[
 E^{N(i)},\qquad h(x,y)=\sum_j\bar x_jy_j,
 \qquad L_{0,i}=\mathcal O^{N(i)}.
                                                               \tag{SF.3}
\]
The space $X'$ is the same construction in dimension $N(m)-1$.
The designated frames are the coordinate bases, not bases obtained by
an existence theorem. Their pairings are $\delta_{ij}$ and their
integer spans are the displayed lattices. Consequently
\texttt{StdTied} holds, $\texttt{gramP}=I$,
$\texttt{gramN}(i)=I$, and membership in
$\texttt{selfDualLatticeCoords}(i)$ is coordinatewise integrality.
In these literal coordinates one can choose $B=I$ in (SF.2).

The \texttt{Pins.iX} clause sends the $j$th basis vector of $X'$ to
the $j$th basis vector of $X_m$. By expanding a vector in its basis,
the inclusion is therefore
\[
 (x_1,\ldots,x_{N(m)-1})\longmapsto
 (x_1,\ldots,x_{N(m)-1},0).                              \tag{SF.4}
\]
There is no change of the underlying class lattices. By definition,
$\texttt{latAt}(u,a,i,c)$ is the original representative lattice pulled
back through the inverse coordinate map. Mapping forward gives exactly
the original representative lattice. The frame set is identified by the
identity map, so its determinant functions are unchanged. These observations
prove all the coordinate assertions of
\texttt{spaces\_in\_coordinates}, under its literal
\texttt{Ok, Pins, Laws, Later} premises.

\paragraph{SF3: the companion identity is a one-coordinate calculation.}
For $0\ne\pi\in\mathcal O$ and $n\ge1$, put
\[
 L_\pi=\mathcal O e_1\oplus\cdots\oplus\mathcal O e_{n-1}
                         \oplus\pi\mathcal O e_n.
\]
Testing (SF.1) against $e_1,\ldots,e_{n-1},\pi e_n$ gives
\[
 L_\pi^\#=
 \mathcal O e_1\oplus\cdots\oplus\mathcal O e_{n-1}
                         \oplus\bar\pi^{-1}\mathcal O e_n.
\]
The same test also proves sufficiency. Hence
\[
 \boxed{\ L_\pi+\bar\pi L_\pi^\#=\mathcal O^n.\ }       \tag{SF.5}
\]
For an explicit decomposition of an integer vector $x$, take
$x'=x-x_ne_n\in L_\pi$ and
$y=(x_n/\bar\pi)e_n\in L_\pi^\#$; then
$x=x'+\bar\pi y$. Conversely each summand in (SF.5) is integral.
The case $n=0$ is the zero module and needs no nonzero-$\pi$ argument.
If $\pi\bar\pi=T\in\mathbb Z$, the same coordinate descriptions give
\[
 L_\pi\subset\mathcal O^n\subset L_\pi^\#,
 \qquad T L_\pi^\#\subset L_\pi.                       \tag{SF.6}
\]
The conjugate on $\pi$ in (SF.5) is required by the convention
that the first argument of $h$ is conjugate linear.

\paragraph{SF4: identify the actual carried lattices.}
The source defines $\texttt{baseL}(n,\pi)$ by
\[
 x_j\in\mathcal O\quad\text{for every }j,
 \qquad x_n\in\pi\mathcal O,
                                                               \tag{SF.7}
\]
using the equivalent zero-indexed condition $j+1=n$ for the last
coordinate. Thus it is precisely $L_\pi$.

Here are the actual joins, rather than an assumption that the companion
formula happens to hold. Supply $D=I.\texttt{rec7}(l,Z)$ and standard
forms $F_u$, Hermitian columns $H$, and level-prime data $A_u$, with
\texttt{StdF}, \texttt{StdLvC}, \texttt{CarH} and \texttt{CarT}
as in the source. At level zero, \texttt{StdF.ends} supplies
\texttt{Reads.lb}, which identifies the forms-side lattice with
\texttt{baseL}(N(0),X.\texttt{piT}(u)). The \texttt{CarH.car0}
identity pulls it back along the exact rank reindexing. Hence the
Hermitian level-zero lattice is $L_{\pi_u}$, with
$\pi_u=X.\texttt{piT}(u)$.

At $1\le i<m$ use \texttt{StdF.mid}, \texttt{Reads.lb} and
\texttt{CarH.carC}. At $i=m$ use \texttt{StdF.ends} and
\texttt{CarH.carM}; the terminal chart set is a singleton.
These give the same description for every chart in its own dimension.
The equalities \texttt{CarT.piU} and \texttt{StdF.piU} identify
the multiplier in every case with the same $\pi_u$.
It is nonzero because $\pi_u\bar\pi_u=T_u$ and $T_u$ is prime.
SF2 identifies the Gram and self-dual lattices. Applying (SF.5) now
proves the complete lattice assertion
\texttt{SymmetricBaseCompatibility}, the chart-indexed version of
\texttt{BaseTied}, including its higher-level clauses. In particular, with
$\mathfrak a_u=(\bar\pi_u)$,
\[
 L_u+\mathfrak a_u L_u^\#=L_0                         \tag{SF.8}
\]
is the actual level-zero \texttt{Companion} identity, for the same
columns used in the comparison. No basis change depending on $u$ has
been introduced.

\paragraph{SF5: a localization lemma with an explicit inverse.}
Let $v$ be a split finite place of the real subfield $F$,
$K=F_v$, $R=\mathcal O_{F_v}$, and let
$j:E\to K$ be the chosen first projection of
$E\otimes_FK\simeq K\times K$ from LI5.
Suppose $L\subset B\mathcal O^n$ and
\[
 N B e_i\in L\quad(1\le i\le n),\qquad N\in R^\times,
                                                               \tag{SF.9}
\]
for an integer $N$. Put $M_v=R\text{-span}\,j(L)$.
Then
\[
 M_v=j(B)R^n,\qquad
 t\in M_v\Longleftrightarrow j(B)^{-1}t\in R^n.         \tag{SF.10}
\]
For the forward inclusion, $B^{-1}x$ is integral for $x\in L$,
and $j(\mathcal O)\subset R$. For the reverse, if
$u=j(B)^{-1}t\in R^n$, then
\[
 t=\sum_i (u_i/N)\,j(NBe_i).                           \tag{SF.11}
\]
Every coefficient belongs to $R$ and every generator comes from $L$.
LI5 identifies $M_v$ with the literal first projection of
$\texttt{locL}_vL$. One can also lift (SF.11) inside that lattice:
the tensor scalars $1\otimes(u_i/N)$ belong to its integral tensor
subring and multiply the images of $NBe_i$. Thus (SF.10) proves the
source's exact existential \texttt{fr\_lat} clause, not just equality
of abstract lattices.

\paragraph{SF6: the same $B$ works at every admissible place.}
For a supplied column level $L=L_u$, integrality of its pairing gives
$L\subset L^\#$, and its \texttt{Lb\_dual} condition gives
$T_u L^\#\subset L$; this extends from \texttt{dualSet} to its span
by linearity. From (SF.8),
\[
 L\subset L_0=B\mathcal O^n\subset L^\#,
 \qquad T_uBe_i\in L.                                  \tag{SF.12}
\]
At a prime admitted by \texttt{GIdx}, $p\nmid T_u$. For every
$v\mid p$, $T_u$ is a unit of $R$: integers $a,b$ with
$aT_u+bp=1$ show that its residue is nonzero. Apply SF5 with $N=T_u$.
Replace only the frame in the already constructed local identification by
\[
 B_v=j_v(B).                                             \tag{SF.13}
\]
Its lattice axiom follows from (SF.10); all other local-identification
fields are unchanged. In particular the completion, splitting map,
prime bijection, embedding $\texttt{iw}$, original orbit representatives
and position map from LI1--LI7 are identical. The alignment identity
with $c^s$ at odd positions and the prescribed-place valuation identity
therefore remain literally the same. This proves
\texttt{LocId.withFr}, \texttt{locIdStd},
\texttt{locIdStd\_fr} and \texttt{locIdStd\_iw}.

\paragraph{SF7: the common-frame conclusion and its boundary.}
Fix one supplied true-column construction and the standard forms and
carrier identifications in SF4. SF2 supplies its standard self-dual
spaces, SF4 supplies its companion identities, and SF1 chooses $B$
once from the common level-zero space. For every word $u$ and every
\texttt{GIdx} place, LI supplies a local identification and SF6 replaces
its frame by (SF.13). This gives simultaneously
\[
 B^*J_0B=I,\qquad
 \pi_{1,v}(\operatorname{locL}_v L_u)=j_v(B)R_v^n
 \quad\text{for all such }u,v.                          \tag{SF.14}
\]
For the actual positive-space coordinates $B=I$ is available. The source
uses the more flexible existence formulation, and both choices satisfy
the same exact matrix and integrality conditions.

This discharges the common-frame \emph{algebra}; it does not construct
the entire supplied standard-column package. Its retained source supplier
is \texttt{colsTiedSym\_std}: besides the lattice identities proved here,
that theorem constructs the forms, levels, characters, residues, class
maps and their simultaneous compatibilities. Its hypotheses
\texttt{Ok, Pins, Laws, Later}, and any earlier host/point existence
needed to supply them, keep their stated status. The class-rank equality,
global comparison-lattice span, spectral/metric estimates and arithmetic
windows require their respective arguments; (SF.14) alone proves none
of those conclusions.
